\subsection{理想}

命题 6. 设 \(\mathrm{R}_1, \ldots, \mathrm{R}_p\) 是 \(\mathbb{R}\) 的不可约分支。对 \(i = 1, \ldots, p\) ，置 \(\mathfrak{g}_i = \mathfrak{h}_{\mathrm{R}_i} + \mathfrak{g}^{\mathrm{R}_i}\) 。则 \(\mathfrak{g}_1, \ldots, \mathfrak{g}_p\) 是 \(\mathfrak{g}\) 的单分量。

诸 \(g_{i}\) 是 g 的理想（命题 1 (ii)）。显然 g 是诸 \(g_{i}\) 的直和，因而是诸 \(g_{i}\) 的积。设 a 与 b 是 g 的互补理想。则 a 与 b 都是半单的且在 ad h 下稳定，故存在 R 的子集 P, Q 使 \(a = h_{P} + g^{P}\) ，\(b = h_{Q} + g^{Q}\) 。于是 \(h_{P}\) ，\(h_{Q}\) 关于 Killing 型在 h 中互为正交补，故 P 与 Q 是 R 的不可约分支的并。这证明了诸 \(g_{i}\) 是 g 的极小理想。

推论 1. g 单当且仅当 R 不可约（换言之，它的邓金图是连通的）。

这由命题 6 得出。

李代数 \(\mathfrak{a}\) 称为绝对单的，如果对每个扩张 \(k'\) （\(k\) 的），\(k'\) -李代数 \(\mathfrak{a}_{(k')}\) 都是单的。

推论 2. 可分裂单李代数是绝对单的。

这由推论 1 得出。

若 g 是 \(A_{l}\) 型（ \(l \geq 1\) ）、\(B_{l}\) 型（ \(l \geq 1\) ）、\(C_{l}\) 型（ \(l \geq 1\) ）或 \(D_{l}\) 型（ \(l \geq 3\) ）的，则称 g 为经典可分裂单李代数。若 g 是 \(E_{6}\) ，\(E_{7}\) ，\(E_{8}\) ，\(F_{4}\) 或 \(G_{2}\) 型的，则称 g 为例外型可分裂单李代数。

